% Intended LaTeX compiler: pdflatex
\documentclass[a4paper,12pt]{article}
\usepackage[utf8]{inputenc}
\usepackage[T1]{fontenc}
\usepackage{graphicx}
\usepackage{longtable}
\usepackage{wrapfig}
\usepackage{rotating}
\usepackage[normalem]{ulem}
\usepackage{amsmath}
\usepackage{amssymb}
\usepackage{capt-of}
\usepackage{hyperref}
\usepackage[margin=0.2in]{geometry}
\usepackage{amsmath}
\usepackage{amsfonts}
\usepackage{indentfirst}
\usepackage{pgffor}
\usepackage{amssymb}
\usepackage{pifont}
\usepackage{cancel}
\usepackage{amsthm}
\usepackage{polynom}
\usepackage{tikz}
\usepackage[tikz]{bclogo}
\usepackage{listings}
\usepackage[most]{tcolorbox}
\usepackage{forest}
\usepackage{adjustbox}
\usepackage{tikz-3dplot}
\usepackage{mathtools}
\usepackage{pgfplots}
\usetikzlibrary{tikzmark,calc,fit,matrix,arrows,automata,positioning}
\usepackage{centernot}
\usepackage{lmodern}
\usepackage{physics}
\usepackage[boxed,linesnumbered,vlined]{algorithm2e}
\usepackage{algpseudocode}
\usepackage{algorithmicx}
\usepackage{svg}
\tikzstyle{every picture}+=[remember picture]
\DeclareMathOperator{\spn}{span}
\DeclareMathOperator{\dom}{domain}
\DeclareMathOperator{\ran}{range}
\DeclareMathOperator{\rng}{range}
\DeclareMathOperator{\img}{Im}
\DeclareMathOperator{\adj}{adj}
\DeclareMathOperator{\Var}{Var}
\DeclareMathOperator{\cov}{cov}
\newcommand\dif[1]{\:\textrm{d}#1}
\DeclarePairedDelimiter\ceil{\lceil}{\rceil}
\DeclarePairedDelimiter\floor{\lfloor}{\rfloor}
\newcommand{\ihat}{\hat{\textbf{\i}}}
\newcommand{\jhat}{\hat{\textbf{\j}}}
\newcommand{\khat}{\hat{\textbf{k}}}
\newcommand{\rhat}{\hat{\textbf{r}}}
\newcommand{\thetahat}{\boldsymbol{\hat{\theta}}}
\newcommand{\?}{\stackrel{?}{=}}
\newcommand{\cmark}{\text{\ding{51}}}
\newcommand{\xmark}{\text{\ding{55}}}
\author{mahmood}
\date{\today}
\title{algorithms sample exam 1}
\hypersetup{
 pdfauthor={mahmood},
 pdftitle={algorithms sample exam 1},
 pdfkeywords={},
 pdfsubject={},
 pdfcreator={Emacs 30.0.50 (Org mode 9.6)}, 
 pdflang={English}}
\begin{document}

\maketitle
\definecolor{Brown}{cmyk}{0,0.81,1,0.60}
\definecolor{OliveGreen}{cmyk}{0.64,0,0.95,0.40}
\definecolor{CadetBlue}{cmyk}{0.62,0.57,0.23,0}
\definecolor{lightlightgray}{gray}{0.9}
\lstset{
  language=C,                             % Code langugage
  basicstyle=\ttfamily,                   % Code font, Examples: \footnotesize, \ttfamily
  keywordstyle=\color{OliveGreen},        % Keywords font ('*' = uppercase)
  commentstyle=\color{gray},              % Comments font
  numbers=left,                           % Line nums position
  numberstyle=\tiny,                      % Line-numbers fonts
  stepnumber=1,                           % Step between two line-numbers
  numbersep=5pt,                          % How far are line-numbers from code
  backgroundcolor=\color{lightlightgray}, % Choose background color
  frame=single,                           % A frame around the code
  tabsize=2,                              % Default tab size
  captionpos=b,                           % Caption-position = bottom
  breaklines=true,                        % Automatic line breaking?
  breakatwhitespace=false,                % Automatic breaks only at whitespace?
  showspaces=false,                       % Dont make spaces visible
  showtabs=false,                         % Dont make tabls visible
  columns=flexible,                       % Column format
  morekeywords={__global__, __device__},  % CUDA specific keywords
  showstringspaces=false,                 % dont show spaces in strings
}
% so that latex doesnt complain about sage not being a language
\lstdefinelanguage{sage}{}

% environment definition for special blocks
\theoremstyle{definition}
\newtheorem{my_example}{Example}
\counterwithin{my_example}{section}
\counterwithin{my_example}{subsection}
\newtheorem{definition}{Definition}
\counterwithin{definition}{section}
\counterwithin{definition}{subsection}
\newtheorem{characteristic}{Characteristic}
\newtheorem{entailment}{Entailment}
%\newtheore*{proof}{Proof}
\newtheorem{lemma}{Lemma}
\newtheorem{question}{Question}
\newtheorem{subquestion}{Subquestion}[question]
%\newtheorem{note}{Note}
\newtheorem{answer}{Answer}
\counterwithin{answer}{question}
\counterwithin{answer}{subquestion}
\newtheorem{step}{Step}[answer]

% the following snippet auto resizes images to fit properly
% Determine if the image is too wide for the page.
% \makeatletter
% \def\ScaleIfNeeded{%
%   \ifdim\Gin@nat@width>\linewidth
%     \linewidth
%   \else
%     \Gin@nat@width
%   \fi
% }
% \makeatother
% % Resize figures that are too wide for the page.
% \let\oldincludegraphics\includegraphics
% \renewcommand\includegraphics[2][]{%
%   \oldincludegraphics[width=\ScaleIfNeeded]{#2}
% }

\newenvironment{note}
  {\begin{bclogo}[logo=\bcinfo, couleurBarre=orange, couleur=lightgray]{ note}}
  {\end{bclogo}}
\newenvironment{attention}
  {\begin{bclogo}[logo=\bcattention, couleurBarre=red, noborder=true, 
                couleur=red!15]{Important!}}
  {\end{bclogo}}

\begin{question}
about \href{20230120221836-select_algorithm.org}{select algorithm}\\[0pt]
\begin{subquestion}
present the \href{20221223133700-computational_problem.org}{computational problem} that the \href{20220706211958-algorithm.org}{algorithm} solves\\[0pt]
\end{subquestion}
\begin{subquestion}
show the \href{20221104220603-algorithm_correctness.org}{correctness} of the algorithm\\[0pt]
\end{subquestion}
\end{question}
\begin{question}
given the following algorithm that receives as input two arrays of characters, the array \(A\) which contains \(n\) elements and the array \(P\) that contains \(m\) elements, such that \(m<n\)\\[0pt]
\includesvg{/Users/mahmooz/brain/out/uK2mk3i}\\[0pt]

\begin{subquestion}
show the computational problem that the algorithm solves\\[0pt]
\begin{answer}
i have no fricking idea, \(P[0]\) is getting dereferenced anyway which is an error\\[0pt]
\end{answer}
\end{subquestion}
\label{orgffda0d4}
\end{question}
\begin{question}
given an array \(A\) whose elements are integers and whose dimensions are \(n \times n\), you have to propose an algorithm that computes the sequence of number in the array whose sum is maximal, the sequence starts with \(A[1,1]\) and ends with \(A[n,n]\), every element in the sequence is one to the right or one below the previous element\\[0pt]
\begin{answer}
we use an array \(B\) in which \(B[i,j]\) is defined as the maximal sum under the constraint that \(B[i,j]\) is in the solution\\[0pt]
\includesvg{/Users/mahmooz/brain/out/oeXKdvm}\\[0pt]

the \href{20221130014441-time_complexity.org}{time complexity} of algorithm is \(\Theta(n)\)\\[0pt]
\end{answer}
\label{org8e7ba57}
\end{question}
\end{document}